\documentclass[11pt]{article}
\usepackage{amsmath}
\usepackage{graphicx}

\title{Keplerian Orbits}
\author{Your Name}
\date{}

\begin{document}

\maketitle

% ============================================================
\section*{Introduction}
% ============================================================

% TODO

% ============================================================
\section*{1. The Gravitational Model}
% ============================================================

% TODO

% ============================================================
\section*{2. The Analytical Orbit}
% ============================================================

% TODO

% ============================================================
\section*{3. From the Equation to a Plot}
% ============================================================

% TODO

% ============================================================
\section*{4. Exploring Eccentricity}
% ============================================================

% TODO

% ============================================================
\section*{5. The Same Problem as a Dynamical System}
% ============================================================

Equation (1) is a second-order system.
With the state vector $\mathbf X = (x, y, v_x, v_y)$ it becomes
\begin{equation}
  \dot{\mathbf X} = \left(v_x,\; v_y,\; -\frac{\mu x}{r^3},\; -\frac{\mu y}{r^3}\right),
\end{equation}
a first-order system that a numerical method can advance step by step.

% ============================================================
\section*{6. Generating the Orbit with Python}
% ============================================================

The script \texttt{generate\_kepler\_data.py} integrates (3) with the
classical fourth-order Runge--Kutta method.
It starts at the periapsis of the orbit in Figure~\ref{fig:analytical},
\begin{equation}
  x_0 = a(1-e) = 0.4, \qquad v_{y0} = \sqrt{\frac{\mu(1+e)}{a(1-e)}} = 2,
\end{equation}
and takes 720 steps over one period, $T = 2\pi\sqrt{a^3/\mu} = 2\pi$.
It writes the trajectory to \texttt{kepler\_orbit.csv}.
The script uses only the Python standard library.

% ============================================================
\section*{7. Plotting the CSV Data}
% ============================================================

% TODO

% ============================================================
\section*{8. Analytical vs.\ Numerical Orbit}
% ============================================================

% TODO

% ============================================================
\section*{9. Kepler's Second Law}
% ============================================================

% TODO

% ============================================================
\section*{10. Kepler's Third Law}
% ============================================================

% TODO

% ============================================================
\section*{11. A Peano--Baker Perspective}
% ============================================================

% TODO

% ============================================================
\section*{12. Peano--Baker and the Kepler Problem}
% ============================================================

% TODO

% ============================================================
\section*{13. From Equations to Computation}
% ============================================================

% TODO

% ============================================================
\section*{Summary}
% ============================================================

% TODO

\end{document}
